\documentclass[11pt]{article}

\usepackage[a4paper,margin=1in]{geometry}
\usepackage{amsmath,amssymb}
\usepackage{hyperref}

\setlength{\parindent}{0pt}
\setlength{\parskip}{0.7em}

\title{Dephasing Diagnostics from Bloch-Vector Length}
\author{}
\date{}

\begin{document}
\maketitle

\section*{Purpose}

Use this note whenever implementing or modifying diagnostics intended to visualize dephasing in custom MCWF results.

For any averaging over different \(N_J\) sectors, group-resolved sectors, or trajectory ensembles, follow \texttt{docs/instructions/bloch\_vector\_averaging.tex}. This file only adds the dephasing-specific rule that the averaged vector length is computed after the relevant vector components have been averaged.

\section*{Goal}

Visualize dephasing by plotting the length of the ensemble-averaged effective \(S\)-Bloch vector:

\begin{equation}
\left|\langle \mathbf S(t)\rangle_{\rm ens}\right|
=
\sqrt{
\langle S_x(t)\rangle_{\rm ens}^2+
\langle S_y(t)\rangle_{\rm ens}^2+
\langle S_z(t)\rangle_{\rm ens}^2
}.
\end{equation}

Dephasing is loss of ensemble phase coherence. In a collective-spin picture, it appears as shrinkage of this averaged vector.

\section*{Core Averaging Rule}

Always average the vector components first, then compute the length:

\begin{equation}
\left|\langle \mathbf S(t)\rangle_{\rm ens}\right|
=
\sqrt{
\langle S_x(t)\rangle_{\rm ens}^2+
\langle S_y(t)\rangle_{\rm ens}^2+
\langle S_z(t)\rangle_{\rm ens}^2
}.
\end{equation}

Do not average per-trajectory lengths:

\begin{equation}
\frac{1}{N_{\rm traj}}
\sum_r
\left|\mathbf S_r(t)\right|.
\end{equation}

The latter measures the typical conditioned-trajectory spin length, not ensemble dephasing.

\section*{Recommended API}

Implement this as standalone post-processing, for example:

\begin{verbatim}
plot_dephasing_bloch_lengths(
    result,
    *,
    normalize=True,
    axes=None,
    output_path=None,
)
\end{verbatim}

The function should accept both \texttt{TrajectoryResult} and \texttt{TrajectoryEnsemble}.

For a single trajectory, compute the vector length from that trajectory's saved observables. For an ensemble, average \(S_x,S_y,S_z\) over trajectories on the shared \texttt{t\_eval} grid before computing the length.

Do not modify MCWF propagation unless the required observables are genuinely unavailable.

\section*{Homogeneous Results}

For homogeneous simulations, plot one curve:

\begin{equation}
\left|\langle \mathbf S_{\rm total}(t)\rangle\right|.
\end{equation}

Use the existing effective-\(S\) observable or moment-extraction path if one exists. Do not invent a second convention for \(S_x,S_y,S_z\).

If effective-\(S\) observables are not available, add only the minimal support needed to compute and post-process them.

\section*{Inhomogeneous Results}

For two-group inhomogeneous simulations, plot one figure with three curves: total, group 1, and group 2.

For groups \(a=1,2\),

\begin{equation}
\left|\langle \mathbf S_a(t)\rangle\right|
=
\sqrt{
\langle S_{x,a}(t)\rangle^2+
\langle S_{y,a}(t)\rangle^2+
\langle S_{z,a}(t)\rangle^2
}.
\end{equation}

For the total curve, sum components first:

\begin{equation}
S_{x,\rm total}=S_{x,1}+S_{x,2},
\qquad
S_{y,\rm total}=S_{y,1}+S_{y,2},
\qquad
S_{z,\rm total}=S_{z,1}+S_{z,2}.
\end{equation}

Then compute

\begin{equation}
\left|\langle \mathbf S_{\rm total}(t)\rangle\right|
=
\sqrt{
\langle S_{x,\rm total}(t)\rangle^2+
\langle S_{y,\rm total}(t)\rangle^2+
\langle S_{z,\rm total}(t)\rangle^2
}.
\end{equation}

If group-resolved effective-\(S\) observables are unavailable, add the minimal support needed to compute them. Keep the diagnostic standalone.

\section*{Normalization}

If \texttt{normalize=True}, divide by the appropriate maximum or reference spin length used elsewhere in the code.

Prefer existing normalization conventions. If no convention exists, use

\begin{equation}
S_{\max}=\frac{N}{2}
\end{equation}

for total homogeneous or effective-spin length, and

\begin{equation}
S_{{\max},a}=\frac{N_a}{2}
\end{equation}

for group-resolved inhomogeneous curves. Document any different active/effective population convention explicitly.

\section*{Optional Relative-Phase Diagnostic}

If group angles are already available, it is useful to also plot

\begin{equation}
\Delta\phi(t)=\phi_1(t)-\phi_2(t).
\end{equation}

This can help distinguish relative dephasing between groups from dephasing within each group.

\section*{Interpretation}

\begin{itemize}
\item If the total length shrinks, ensemble coherence is being lost.
\item If the total length shrinks while group-resolved lengths remain large, the dominant effect is likely relative dephasing between groups.
\item If group-resolved lengths also shrink, there is dephasing within the groups as well.
\end{itemize}

\end{document}
